\clearpage\thispagestyle{empty}
\pdfbookmark[0]{Quellen zu den geozentrischen Positionen der Planeten (1804)}{bemerkungen1804-sources}
\begin{center}\textbf{Quellen zu den geozentrischen Positionen der Planeten (1804)}\end{center}
\noindent Diese KI-Übertragung in heutiges Deutsch folgt Carl Friedrich Gauß’ deutschem Original in den \textit{Werken}, Band~VI (Göttingen, 1874), S.~94--105, einschließlich der am Ende abgedruckten handschriftlichen Bemerkungen. Der Aufsatz erschien zuerst als Beitrag L im Maiheft 1804 der von Franz Xaver von Zach herausgegebenen \textit{Monatlichen Correspondenz zur Beförderung der Erd- und Himmels-Kunde}, Band~9, S.~385--400. \href{https://gdz.sub.uni-goettingen.de/id/PPN236007467}{Historischer deutscher Text in den Werken}. \href{https://www.digitale-sammlungen.de/en/view/bsb10538601?page=421}{Erstdruck}.
\par\smallskip\noindent In den \textit{Werken} steht \(x-X\) im Nenner der Formel für \(\operatorname{tang}\alpha\) (S.~96). Im Erstdruck steht dort \(x-Z\) (S.~387), ebenso im Nachdruck in Anhang VII, S.~53--58, zu \href{https://gdz.sub.uni-goettingen.de/id/PPN60187157X}{Carl Haases deutscher Übersetzung der \textit{Theoria motus}} (Hannover, 1865). In den handschriftlichen Bemerkungen der \textit{Werke} bedeutet \textit{in plano ipsius} \(Z\) „in der Ebene von \(Z\)“.
\par\smallskip\noindent Gauß behandelt das verwandte Problem in Artikel~53--57, S.~60--66, der \href{https://www2.karlin.mff.cuni.cz/~tuma/Aplikace21/theoryofmotionof00gausuoft.pdf#page=89}{\textit{Theoria motus} in Charles Henry Davis’ englischer Übersetzung} (Boston, 1857). Artikel~56, S.~65, verweist für eine weitere Lösung auf den Aufsatz vom Mai 1804.
\par\smallskip\noindent Übertragung und Satz wurden mit KI erstellt.
\clearpage
